\documentclass[a4paper,fontset=none]{ctexart}

\usepackage{anyfontsize}
\usepackage{amsmath,amssymb,mathrsfs,wasysym}
\usepackage{autobreak}
\allowdisplaybreaks
\usepackage[T1]{fontenc}
\usepackage{textcomp}
\usepackage{amsthm}
\usepackage[libertine,vvarbb]{newtxmath}
\usepackage[scr=rsfso,frak=euler,bb=ams]{mathalfa}
\usepackage{bm}
\usepackage{indentfirst}
\setlength{\parindent}{0em}
\usepackage{parskip}
\usepackage{geometry}
\geometry{top=3cm,bottom=3cm,left=2cm,right=2cm}
\usepackage{fontspec}
\setmainfont{Linux Libertine O}
\setsansfont{Linux Biolinum O}
\setmonofont{JetBrains Mono}
\setCJKmainfont{SimSun}[BoldFont=Noto Sans CJK SC Medium]

\begin{document}

    \title{\textbf{天体物理导论作业}}
    \author{1900011604\ \ 张植竣\ \ 北京大学}
    \date{2021年4月20日}
    \maketitle

    \section*{第六章}

    \begin{itemize}

        \item[1.]
        推导$r_{\mathrm{co}}$如下：
        
        由Kepler第三定律得：
        \[
        \frac{r_{\mathrm{co}}^3}{P^2} = \frac{GM}{4\pi^2}
        \]
        整理得：
        \[
        r_{\mathrm{co}} = \left(\frac{GM}{4\pi^2}\right)^{\tfrac{1}{3}}P^{\tfrac{2}{3}}
        \]
        推导$r_{\mathrm{m}}$如下：
        
        由于磁偶极子产生磁场的磁感应强度分布$B \sim \mu r^{-3}$，取$\mu = BR^3$，其中$B$为中子星表面磁场，则$r = r_{\mathrm{m}}$处磁感应强度为：
        \[
        B_{r_{\mathrm{m}}} = \frac{BR^3}{r_{\mathrm{m}}^3}
        \]
        由粒子的动能密度与磁能密度相等得：
        \[
        \frac{B_{r_{\mathrm{m}}}^2}{8\pi} = \frac{1}{2}\rho v^2
        \]
        假设粒子在到达$r_{\mathrm{m}}$前近似自由落体，则有：
        \[
        v = \sqrt{\frac{2GM}{r_{\mathrm{m}}^2}}
        \]
        又因为：
        \[
        \dot{M} = 4\pi r_{\mathrm{m}}^2\rho v
        \]
        则有：
        \[
        \frac{B_{r_{\mathrm{m}}}^2}{8\pi} = \frac{1}{2}v^2\cdot\frac{\dot{M}}{4\pi r_{\mathrm{m}}^2 v} = \frac{\dot{M}v}{8\pi r_{\mathrm{m}}^2}
        \]
        即：
        \[
        B_{r_{\mathrm{m}}}^2 = \frac{\dot{M}\sqrt{2GM}}{r_{\mathrm{m}}^2\sqrt{r_{\mathrm{m}}}}
        \]
        整理得：
        \[
        r_{\mathrm{m}}^{-\tfrac{5}{2}} = \frac{B_{r_{\mathrm{m}}}^2}{\dot{M}\sqrt{2GM}}
        \]
        代入$B_{r_{\mathrm{m}}}$，即得：
        \[
        r_{\mathrm{m}}^{\tfrac{7}{2}} = \frac{R^6 B^2}{\dot{M}\sqrt{2GM}}
        \]
        故磁层半径为：
        \[
        r_{\mathrm{m}} = \left({\frac{R^6 B^2}{\dot{M}\sqrt{2GM}}}\right)^{\tfrac{2}{7}}
        \]

        \item[3.]
        大规模吸积发生的必要条件一般为：$R<r_{\mathrm{m}}<r_{\mathrm{co}}<r_{\mathrm{L}}$，具体上需要满足：$r_{\mathrm{m}}<r_{\mathrm{br}}<r_{\mathrm{co}}$，即：
        \[
        \omega_{\mathrm{s}} = \left(\frac{r_{\mathrm{m}}}{r_{\mathrm{co}}}\right)^{\tfrac{3}{2}} < \omega_{\mathrm{c}} = 0.35
        \]
        其中：
        \[
        r_{\mathrm{m}} = \left({\frac{R^6 B^2}{\dot{M}\sqrt{2GM}}}\right)^{\tfrac{2}{7}},\qquad
        \dot{M} = \pi R_{\mathrm{a}}^2 V\rho,\qquad
        R_{\mathrm{a}} = \frac{GM}{V^2}
        \]
        代入数据可以计算出：
        \[
        r_{\mathrm{m}} = 1.4310\times10^{11}\,\mathrm{cm}
        \]
        而：
        \[
        r_{\mathrm{co}} = \left(\frac{GM}{4\pi^2}\right)^{\tfrac{1}{3}}P^{\tfrac{2}{3}} = CP^{\tfrac{2}{3}}
        \]
        代入数据计算得：
        \[
        C = 1.8337\times10^{12}\,\mathrm{cm \cdot s^{-\tfrac{1}{3}}}
        \]
        则有：
        \[
        P > \frac{r_{\mathrm{m}}^{\tfrac{3}{2}}}{\omega_{\mathrm{c}}C^{\tfrac{3}{2}}} = 8.4345\times10^{4}\,\mathrm{s}
        \]

        \item[\textbf{补充题}]
        不是。因为暗物质粒子之间的相互作用极弱，在下落过程中无法达到热平衡，因而无法定义温度；而Bondi吸积需要气体达到热平衡。实际上暗物质粒子的下落过程可以近似用双曲轨道描述，只有靠近黑洞视界的粒子才可能被吸积。对于Schwarzschild黑洞可以用Schwarzschild半径作为吸积半径，Kerr黑洞的吸积率更加复杂，但都不是Bondi吸积率。

    \end{itemize}

    \section*{第七章}

    \begin{itemize}

        \item[\textbf{补充题}]
        推导极端相对论性理想气体状态方程如下：

        在气体中任取一面元$\Delta S$，并取$x$方向沿其法线方向。考虑速度为$\bm{v_i}$的粒子，设其数密度为$n_i$。在$\Delta t$时间内，只有处于以$\Delta S$为底，以$\bm{v_i}\cdot{\bm{x}}\Delta t$为高的柱体内的粒子才能与$\Delta S$相碰。这里显然要求$\bm{v_i}\cdot\bm{x} > 0$，记该组粒子的数密度为$n_i^+$，则该柱体中的粒子数为：
        \[
        N_i = n_i^+\Delta S(\bm{v_i}\cdot{\bm{x}})\Delta t
        \]
        传递的动量为：
        \[
        \Delta\bm{p_i} = N_i\cdot(2\bm{p_i}\cdot\bm{x}) = 2n_i^+\Delta S(\bm{v_i}\cdot{\bm{x}})\Delta t(\bm{p_i}\cdot\bm{x})\hat{x}
        \]
        由于气体分子速度分布各向同性，$n_i^+ = n_i^-$，即$2n_i^+ = n_i$，于是有：
        \[
        \Delta\bm{p_i} = n_i\Delta S(\bm{v_i}\cdot{\bm{x}})\Delta t(\bm{p_i}\cdot\bm{x})\hat{x}
        \]
        施加的压强为：
        \[
        P_i = \frac{\Delta\bm{p_i}}{\Delta t\Delta S} = n_i(\bm{v_i}\cdot{\bm{x}})(\bm{p_i}\cdot\bm{x})
        \]
        总压强为：
        \begin{align*}
            P
            &= \sum_i P_i \\
            &= \sum_i n_i(\bm{v_i}\cdot{\bm{x}})(\bm{p_i}\cdot\bm{x}) \\
            &= n\sum_i \frac{n_i(\bm{v_i}\cdot{\bm{x}})(\bm{p_i}\cdot\bm{x})}{n} \\
            &= n\overline{\bm{p_x}\cdot\bm{v_x}}
        \end{align*}
        由于粒子热运动各向同性，$\overline{\bm{p_x}\cdot\bm{v_x}} = \overline{\bm{p_y}\cdot\bm{v_y}} = \overline{\bm{p_z}\cdot\bm{v_z}}$，则$\overline{\bm{p}\cdot\bm{v}} = 3\overline{\bm{p_x}\cdot\bm{v_x}}$，即有：
        \[
        P = \frac{1}{3}n\overline{\bm{p}\cdot\bm{v}}
        \]
        对于极端相对论性理想气体，$v \approx c$，$p \approx mc$，单个粒子的动能$\varepsilon_\mathrm{k} \approx mc^2$，则：
        \[
        P = \frac{1}{3}n\overline{mv^2} = \frac{1}{3}nmc^2 = \frac{1}{3}n\varepsilon_\mathrm{k}
        \]
        记粒子的能量密度$\rho = n\varepsilon_\mathrm{k}$，即有：
        \[
        P = \frac{1}{3}\rho
        \]

    \end{itemize}

\end{document}